\documentclass[12pt]{article}
\usepackage[a4paper, margin=1in]{geometry}
\usepackage{amsmath, amssymb, minted, enumitem, cmbright, hyperref}
\renewcommand{\thesection}{\Alph{section}}
\renewcommand{\thesubsection}{\thesection.\arabic{subsection}}
\renewcommand{\t}{\texttt}
\DeclareMathOperator{\dist}{dist}
\begin{document}
\section{Application Questions}
\textit{[Note: The original questions were written in C++, so I provided a Python translation in the associated \t{.py} file and in Appendix 1. The solutions are based on these Python versions.]}
\subsection{Archery Practice}
\textit{[Note: The original question used C++ \t{std::set} which is different from Python \t{set}. In particular, it maintains sorted order. The Python equivalent would be \t{SortedList}. But actually the optimal solution does not need to use this.]}

\begin{enumerate}
\item It uses the correct data structure that we learn in class, i.e., set \fbox{NO, I will elaborate}
\item Ignoring time complexity, it gives correct output for all valid inputs: \fbox{NO, I will elaborate}
\item Ignoring correctness, it runs in: $\boxed{O(nk)}$
\item It can crash on at least one valid corner case input: \fbox{NO}
\item I can improve upon the Generated AI code: \fbox{YES, I will elaborate}
\end{enumerate}

The original code is not correct because a set removes duplicate distances, but different arrows may have the same distance from the origin. For example, shots $(1,0)$ and $(0,1)$ are different coordinates but both have distance $1$.

The running time is $O(nk)$, since each insertion costs $O(\log n)$* and finding the $k^{\text{th}}$ element costs $O(k)$. A better solution is to use a max-heap, keeping the $k$ smallest distances seen so far. Then the top of the heap is the $k^{\text{th}}$ nearest distance. This gives $O(n\log k)$ time complexity.

* since we are using \t{SortedList}.

\begin{minted}{python}
def archery(self, shots: List[Tuple[int, int]], k: int) -> int:
    from heapq import heappush, heappop
    h = []
    ans = 0
    for x, y in shots:
        d = abs(x) + abs(y)
        heappush(h, -d)
        if len(h) > k:
            heappop(h)
        if len(h) == k:
            ans += -h[0]
    return ans
\end{minted}
\newpage
\subsection{Integer Sibling Groups}
\textit{[Note: The original question used UFDS which is non-examinable in IT5003. But actually the optimal solution does not need to use this.]}

\begin{enumerate}
\item It uses the correct data structure that we learn in class, i.e., UFDS \fbox{YES}
\item Ignoring time complexity, it gives correct output for all valid inputs: \fbox{YES}
\item Ignoring correctness, it runs in: $\boxed{O\left(n^2\right)}$
\item It can crash on at least one valid corner case input: \fbox{NO}
\item I can improve upon the Generated AI code: \fbox{YES, I will elaborate}
\end{enumerate}

UFDS can group numbers that have the same digit sum and every pair of numbers with the same digit sum gets unioned, so each digit-sum group becomes one connected component and the final number of disjoint sets equals the number of sibling groups. As the integers are less than $10^7$, \t{get\_digit\_sum} runs in constant time. There are two nested loops over all pairs of numbers so the time complexity is $O\left(n^2\right)$.

However, UFDS is not necessary. A better solution is to compute the digit sum of each number and insert it into a set. The answer is just the number of distinct digit sums, which runs in $O(n)$ with a set.

\begin{minted}{python}
def countSiblingGroups(self, S: Set[int]) -> int:
    digit_sums = set()
    for x in S:
        digit_sums.add(self.get_digit_sum(x))
    return len(digit_sums)

def get_digit_sum(self, n: int) -> int:
    total = 0
    while n > 0:
        total += n % 10
        n //= 10
    return total
\end{minted}
\newpage
\subsection{VisuAlgo AVL Tree Visualization}
\textit{[Note: The original question used AVL tree which is non-examinable in IT5003. But the code and solution are included for completeness.]}

\begin{itemize}
\item My analysis for SolutionTwin is:
\begin{enumerate}
\item Ignoring time complexity, it gives correct output for all valid inputs: \fbox{YES}
\item Ignoring correctness, it runs in: $\boxed{O\left(n\right)}$
\item It can crash on at least one valid corner case input: \fbox{NO}
\end{enumerate}
\item My analysis for SolutionTalk is:
\begin{enumerate}
\item Ignoring time complexity, it gives correct output for all valid inputs: \fbox{YES}
\item Ignoring correctness, it runs in: $\boxed{O\left(n\right)}$
\item It can crash on at least one valid corner case input: \fbox{NO}
\end{enumerate}
\item Thus, I will use: \fbox{Either solution}
\end{itemize}

Both SolutionTwin and SolutionTalk correctly balance a BST by first performing an inorder traversal to obtain a sorted list of values, and then rebuilding a balanced BST using a divide-and-conquer approach. This guarantees correctness.

Both solutions run in $O(n)$ time, since each node is visited once during traversal and once during reconstruction.
\newpage
\subsection{Strongest Trinity}
\begin{enumerate}
\item It uses the correct graph data structure (check space complexity): \fbox{NO, I will elaborate}
\item Ignoring time complexity, it gives correct output for all valid inputs: \fbox{YES}
\item Ignoring correctness, it runs in: $\boxed{O\left(n^4\right)}$
\item It can crash on at least one valid corner case input: \fbox{NO}
\item I can improve upon the Generated AI code: \fbox{YES, I will elaborate}
\end{enumerate}

The original code correctly identifies all trinities and computes their strength by counting connections to other nodes. Therefore, it produces correct results for all valid inputs.

However, it uses AM with $O\left(n^2\right)$ space, which is inefficient. We should use an AL with $O(n+m)$ space. (In general, we should avoid AM unless the graph is sparse and/or we need to frequently check existence of edges.)

The time complexity is $O\left(n^4\right)$ since there are four nested loops.

This solution can be significantly improved. Note that the strength of a trinity $(u,v,w)$ is simply $\deg(u)+\deg(v)+\deg(w)-6$, i.e., the number of edges outside the trinity minus the number of edges inside the trinity. Instead of checking all trinities, we check every edge $(u, v)$ with their common neighbors, which takes time $O(\min(\deg(u),\deg(v))$. Summing over all pairs of $(u, v)$ gives $O(nm)=O\left(n^3\right)$ time in the worst case*.

* where the graph is dense, i.e., $\deg(v)\in O(n), m\in O\left(n^2\right)$.
\begin{minted}{python}
def maxTrinityStrength(self, n: int, EL: List[Tuple[int, int]]) -> int:
    adj = [set() for _ in range(n + 1)]
    for u, v in EL:
        adj[u].add(v)
        adj[v].add(u)
    deg = [len(adj[i]) for i in range(n + 1)]
    return max(deg[u] + deg[v] + deg[w] for u, v in EL
               for w in adj[u] & adj[v]) - 6
\end{minted}
\newpage
\subsection{Yoshi's Tasks}
We can use Kahn's algorithm with a min-heap. The idea is as follows:
\begin{itemize}
\item Build a directed graph with edge $i\to j$ for each precedence pair $\{i, j\}$.
\item Compute the number of unfinished prerequisites for each task, i.e., the indegree.
\item Insert every task with indegree $0$ into a min-heap.
\item Repeatedly:
\begin{itemize}
\item remove the smallest task $u$ from the heap,
\item append $u$ to the answer,
\item for every edge $u\to v$, decrement the indegree of $v$,
\item if the indegree of $v$ is $0$, insert $v$ into the heap.
\end{itemize}
\item If the answer contains fewer than $n$ tasks, there is a cycle, so return $-1$. Otherwise return the answer.
\end{itemize}

Each task is inserted and removed from the heap once, costing $O(n\log n)$. Each edge is processed once, and heap insertions cost at most $O(m\log n)$ total. Thus in total this is $O((n+m)\log n)$.

\begin{minted}{python}
def yoshiTasks(self, n: int, edges: List[Tuple[int, int]]) -> List[int]:
    from heapq import heapify, heappush, heappop
    g = [[] for _ in range(n + 1)]
    d = [0] * (n + 1)
    for u, v in edges:
        g[u].append(v)
        d[v] += 1
    h = [i for i in range(1, n + 1) if d[i] == 0]
    heapify(h)
    ans = []
    while h:
        u = heappop(h)
        ans.append(u)
        for v in g[u]:
            d[v] -= 1
            if d[v] == 0:
                heappush(h, v)
    return ans if len(ans) == n else [-1]
\end{minted}
\newpage
\subsection{The Last Question}
\textit{[Note: This question is similar to CS2040C 23S1 Final B.3.2 (ASSSP).]}

We can run Dijkstra twice, from $0$ and from $n-1$. An edge $(u, v, w)$ is on some shortest path if and only if going through it can still achieve the total shortest distance. The idea is as follows:
\begin{itemize}
\item Run Dijkstra from $0$ to get an array $\dist_0$ where each $\dist_0[u]$ denotes the shortest distance between $0$ and $u$.
\item Run Dijkstra from $n-1$ to get an array $\dist_{n-1}$ where each $\dist_{n-1}[v]$ denotes the shortest distance between $n-1$ and $v$.
\item For any edge $(u,v,w)$, it is part of at least one shortest path if and only if

$\dist_0[u]+w+\dist_{n-1}[v]=\dist_0[n-1]$ or $\dist_0[v]+w+\dist_{n-1}[u]=\dist_0[n-1]$.

Thus we can just enumerate all edges and check the condition.
\end{itemize}

The time complexity is dominated by Dijkstra which is $O((n+m)\log n)$.

\begin{minted}{python}
def theLastQuestion(self, n: int, edges: List[Tuple[int, int, int]]) -> int:
    from heapq import heappush, heappop
    from math import inf

    def dijkstra(s):
        g = [[] for _ in range(n)]
        for u, v, w in edges:
            g[u].append((v, w))
            g[v].append((u, w))
        dist = [inf] * n
        dist[s] = 0
        q = [(0, s)]
        while q:
            d, u = heappop(q)
            if d != dist[u]:
                continue
            for v, w in g[u]:
                if d + w < dist[v]:
                    dist[v] = d + w
                    heappush(q, (dist[v], v))
        return dist
    d1 = dijkstra(0)
    d2 = dijkstra(n - 1)
    d = d1[n - 1]
    return sum(d1[u] + w + d2[v] == d or d1[v] + w + d2[u] == d
               for u, v, w in edges)
\end{minted}
\newpage
\section*{Appendix 1: Python Code [Original]}
\begin{minted}{python}
from typing import List, Tuple, Optional, Set


class TreeNode:
    def __init__(self, value=0, left=None, right=None):
        self.value = value
        self.left = left
        self.right = right


class UnionFind:
    def __init__(self, n: int):
        self.parent = list(range(n))
        self.rank = [0] * n
        self.num_sets = n

    def find(self, x: int) -> int:
        if self.parent[x] != x:
            self.parent[x] = self.find(self.parent[x])
        return self.parent[x]

    def union_set(self, a: int, b: int) -> None:
        ra, rb = self.find(a), self.find(b)
        if ra == rb:
            return
        if self.rank[ra] < self.rank[rb]:
            ra, rb = rb, ra
        self.parent[rb] = ra
        if self.rank[ra] == self.rank[rb]:
            self.rank[ra] += 1
        self.num_sets -= 1

    def num_disjoint_sets(self) -> int:
        return self.num_sets


class Solution:
    # A.1 Archery Practice
    def archery(self, shots: List[Tuple[int, int]], k: int) -> int:
        from sortedcontainers import SortedList
        S = SortedList()
        ans = 0
        for x, y in shots:
            d = abs(x) + abs(y)
            S.add(d)
            if len(S) >= k:
                ans += sorted(S)[k - 1]
        return ans

    # A.2 Integer Sibling Groups
    def countSiblingGroups(self, S: Set[int]) -> int:
        nums = list(S)
        n = len(nums)
        if n == 0:
            return 0
        uf = UnionFind(n)
        digit_sums = [self.get_digit_sum(x) for x in nums]
        for i in range(n):
            for j in range(i + 1, n):
                if digit_sums[i] == digit_sums[j]:
                    uf.union_set(i, j)
        return uf.num_disjoint_sets()

    def get_digit_sum(self, n: int) -> int:
        total = 0
        while n > 0:
            total += n % 10
            n //= 10
        return total

    # A.3 VisuAlgo AVL Tree Visualization
    # SolutionTwin version
    def balanceBST_twin(self, root: Optional[TreeNode]) -> Optional[TreeNode]:
        nodes: List[int] = []
        self.inorder_twin(root, nodes)
        return self.build_twin(nodes, 0, len(nodes) - 1)

    def inorder_twin(self, root: Optional[TreeNode], nodes: List[int]) -> None:
        if not root:
            return
        self.inorder_twin(root.left, nodes)
        nodes.append(root.value)
        self.inorder_twin(root.right, nodes)

    def build_twin(self, nodes: List[int], l: int, r: int) -> Optional[TreeNode]:
        if l > r:
            return None
        m = (l + r) // 2
        node = TreeNode(nodes[m])
        node.left = self.build_twin(nodes, l, m - 1)
        node.right = self.build_twin(nodes, m + 1, r)
        return node

    # SolutionTalk version
    def __init__(self):
        self.a = []

    def balanceBST_talk(self, root: Optional[TreeNode]) -> Optional[TreeNode]:
        # self.a = []
        self.inorder_talk(root)
        return self.build_talk(0, len(self.a) - 1)

    def inorder_talk(self, root: Optional[TreeNode]) -> None:
        if not root:
            return
        self.inorder_talk(root.left)
        self.a.append(root.value)
        self.inorder_talk(root.right)

    def build_talk(self, l: int, r: int) -> Optional[TreeNode]:
        if l > r:
            return None
        m = (l + r) >> 1
        node = TreeNode(self.a[m])
        node.left = self.build_talk(l, m - 1)
        node.right = self.build_talk(m + 1, r)
        return node

    # A.4 Strongest Trinity
    def maxTrinityStrength(self, n: int, EL: List[Tuple[int, int]]) -> int:
        AM = [[0] * (n + 1) for _ in range(n + 1)]
        for u, v in EL:
            AM[u][v] = AM[v][u] = 1
        ans = -1
        for i in range(1, n + 1):
            for j in range(i + 1, n + 1):
                for k in range(j + 1, n + 1):
                    if AM[i][j] and AM[j][k] and AM[i][k]:
                        strength = 0
                        for x in range(1, n + 1):
                            if x in (i, j, k):
                                continue
                            if AM[i][x]:
                                strength += 1
                            if AM[j][x]:
                                strength += 1
                            if AM[k][x]:
                                strength += 1
                        ans = max(ans, strength)
        return ans


tests = [[[(1, 1), (3, 4), (2, 1), (-1, 0)], 2], [[(1, 0), (0, 1), (2, 0)], 2]]
for t in tests:
    print(f"{t} -> {Solution().archery(*t)}")
"""[[(1, 1), (3, 4), (2, 1), (-1, 0)], 2] -> 12
[[(1, 0), (0, 1), (2, 0)], 2] -> 2"""

tests = [{5, 6, 7, 23, 24, 33, 51, 123}, {12, 21, 30, 5}, {10, 20, 11, 101, 99}]
for t in tests:
    print(f"{t} -> {Solution().countSiblingGroups(t)}")
"""{33, 5, 6, 7, 51, 23, 24, 123} -> 3
{5, 12, 21, 30} -> 2
{99, 20, 101, 10, 11} -> 3"""


def inorder(root):
    if not root:
        return []
    return inorder(root.left) + [root.value] + inorder(root.right)


def height(root):
    if not root:
        return -1
    return 1 + max(height(root.left), height(root.right))


root = TreeNode(37)
root.left = TreeNode(28)
root.left.left = TreeNode(24)
root.right = TreeNode(44)
root.right.right = TreeNode(83)
root.right.right.left = TreeNode(56)
root.right.right.left.left = TreeNode(46)
root.right.right.left.right = TreeNode(63)
root.right.right.right = TreeNode(93)
root.right.right.right.left = TreeNode(91)
b1 = Solution().balanceBST_twin(root)
b2 = Solution().balanceBST_talk(root)
print(f"{inorder(root), height(root)} -> {inorder(b1), height(b1)}")
print(f"{inorder(root), height(root)} -> {inorder(b2), height(b2)}")

tests = [[7, [(1, 2), (1, 3), (2, 3), (1, 4), (2, 5), (3, 6), (6, 7)]],
         [7, [(1, 2), (1, 3), (1, 4), (2, 5), (3, 6), (6, 7)]],
         [6, [(1, 2), (1, 3), (2, 3), (1, 4), (2, 5), (3, 6)]],
         [5, [(1, 2), (2, 3), (1, 3), (1, 4), (2, 4), (3, 4)]]]
for t in tests:
    print(f"{t} -> {Solution().maxTrinityStrength(*t)}")
"""[7, [(1, 2), (1, 3), (2, 3), (1, 4), (2, 5), (3, 6), (6, 7)]] -> 3
[7, [(1, 2), (1, 3), (1, 4), (2, 5), (3, 6), (6, 7)]] -> -1
[6, [(1, 2), (1, 3), (2, 3), (1, 4), (2, 5), (3, 6)]] -> 3
[5, [(1, 2), (2, 3), (1, 3), (1, 4), (2, 4), (3, 4)]] -> 3"""
\end{minted}
\newpage
\section*{Appendix 2: Python Code [Optimal]}
\begin{minted}{python}
from typing import List, Tuple, Set


class Solution:
    # A.1 Archery Practice
    def archery(self, shots: List[Tuple[int, int]], k: int) -> int:
        from heapq import heappush, heappop
        h = []
        ans = 0
        for x, y in shots:
            d = abs(x) + abs(y)
            heappush(h, -d)
            if len(h) > k:
                heappop(h)
            if len(h) == k:
                ans += -h[0]
        return ans

    # A.2 Integer Sibling Groups
    def countSiblingGroups(self, S: Set[int]) -> int:
        digit_sums = set()
        for x in S:
            digit_sums.add(self.get_digit_sum(x))
        return len(digit_sums)

    def get_digit_sum(self, n: int) -> int:
        total = 0
        while n > 0:
            total += n % 10
            n //= 10
        return total

    # A.4 Strongest Trinity
    def maxTrinityStrength(self, n: int, EL: List[Tuple[int, int]]) -> int:
        adj = [set() for _ in range(n + 1)]
        for u, v in EL:
            adj[u].add(v)
            adj[v].add(u)
        deg = [len(adj[i]) for i in range(n + 1)]
        return max((deg[u] + deg[v] + deg[w] - 6 for u, v in EL
                    for w in adj[u] & adj[v]), default=-1)

    # A.5 Yoshi's Tasks
    def yoshiTasks(self, n: int, edges: List[Tuple[int, int]]) -> List[int]:
        from heapq import heapify, heappush, heappop
        g = [[] for _ in range(n + 1)]
        d = [0] * (n + 1)
        for u, v in edges:
            g[u].append(v)
            d[v] += 1
        h = [i for i in range(1, n + 1) if d[i] == 0]
        heapify(h)
        ans = []
        while h:
            u = heappop(h)
            ans.append(u)
            for v in g[u]:
                d[v] -= 1
                if d[v] == 0:
                    heappush(h, v)
        return ans if len(ans) == n else [-1]

    # A.6 The Last Question
    def theLastQuestion(self, n: int,
                        edges: List[Tuple[int, int, int]]) -> int:
        from heapq import heappush, heappop
        from math import inf

        def dijkstra(s):
            g = [[] for _ in range(n)]
            for u, v, w in edges:
                g[u].append((v, w))
                g[v].append((u, w))
            dist = [inf] * n
            dist[s] = 0
            q = [(0, s)]
            while q:
                d, u = heappop(q)
                if d != dist[u]:
                    continue
                for v, w in g[u]:
                    if d + w < dist[v]:
                        dist[v] = d + w
                        heappush(q, (dist[v], v))
            return dist
        d1 = dijkstra(0)
        d2 = dijkstra(n - 1)
        d = d1[n - 1]
        return sum(d1[u] + w + d2[v] == d or d1[v] + w + d2[u] == d
                   for u, v, w in edges)


tests = [[[(1, 1), (3, 4), (2, 1), (-1, 0)], 2], [[(1, 0), (0, 1), (2, 0)], 2]]
for t in tests:
    print(f"{t} -> {Solution().archery(*t)}")
"""[[(1, 1), (3, 4), (2, 1), (-1, 0)], 2] -> 12
[[(1, 0), (0, 1), (2, 0)], 2] -> 2"""

tests = [{5, 6, 7, 23, 24, 33, 51, 123}, {12, 21, 30, 5}, {10, 20, 11, 101, 99}]
for t in tests:
    print(f"{t} -> {Solution().countSiblingGroups(t)}")
"""{33, 5, 6, 7, 51, 23, 24, 123} -> 3
{5, 12, 21, 30} -> 2
{99, 20, 101, 10, 11} -> 3"""

tests = [[7, [(1, 2), (1, 3), (2, 3), (1, 4), (2, 5), (3, 6), (6, 7)]],
         [7, [(1, 2), (1, 3), (1, 4), (2, 5), (3, 6), (6, 7)]],
         [6, [(1, 2), (1, 3), (2, 3), (1, 4), (2, 5), (3, 6)]],
         [5, [(1, 2), (2, 3), (1, 3), (1, 4), (2, 4), (3, 4)]]]
for t in tests:
    print(f"{t} -> {Solution().maxTrinityStrength(*t)}")
"""[7, [(1, 2), (1, 3), (2, 3), (1, 4), (2, 5), (3, 6), (6, 7)]] -> 3
[7, [(1, 2), (1, 3), (1, 4), (2, 5), (3, 6), (6, 7)]] -> -1
[6, [(1, 2), (1, 3), (2, 3), (1, 4), (2, 5), (3, 6)]] -> 3
[5, [(1, 2), (2, 3), (1, 3), (1, 4), (2, 4), (3, 4)]] -> 3"""

tests = [[5, [(1, 4), (4, 2), (1, 2), (1, 5)]],
         [5, [(1, 4), (4, 2), (2, 1), (1, 5)]]]
for t in tests:
    print(f"{t} -> {Solution().yoshiTasks(*t)}")
"""[5, [(1, 4), (4, 2), (1, 2), (1, 5)]] -> [1, 3, 4, 2, 5]
[5, [(1, 4), (4, 2), (2, 1), (1, 5)]] -> [-1]"""

tests = [[6, [(0, 1, 1), (0, 2, 1), (0, 3, 4), (1, 4, 3),
              (2, 3, 2), (3, 4, 1), (4, 5, 3)]]]
for t in tests:
    print(f"{t} -> {Solution().theLastQuestion(*t)}")
"""[6, [(0, 1, 1), (0, 2, 1), (0, 3, 4), (1, 4, 3),
        (2, 3, 2), (3, 4, 1), (4, 5, 3)]] -> 6"""
\end{minted}
\end{document}
